% =====================================================================
%  CARNET D'ENTRAINEMENT — FICHE 11 — THEME 4
%  Tuto : Potentiels standard et prévision des réactions
% =====================================================================
\documentclass[11pt,a4paper]{article}
\usepackage[utf8]{inputenc}
\usepackage[T1]{fontenc}
\usepackage{lmodern}
\usepackage[french]{babel}
\usepackage{amsmath,amssymb}
\usepackage[version=4]{mhchem}
\usepackage{siunitx}
\sisetup{output-decimal-marker={,},per-mode=symbol}
\usepackage{xcolor}
\usepackage{array}
\usepackage{booktabs}
\usepackage{enumitem}
\usepackage[most]{tcolorbox}
\usepackage{tikz}
\usetikzlibrary{arrows.meta}
\usepackage[a4paper,margin=2.0cm,headheight=26pt,headsep=10pt]{geometry}
\usepackage{fancyhdr}
\usepackage{titlesec}
\usepackage[hidelinks]{hyperref}

\definecolor{primary}{RGB}{146,95,160}
\definecolor{primarydark}{RGB}{92,52,104}
\definecolor{defcol}{RGB}{0,114,178}
\definecolor{thmcol}{RGB}{0,140,120}
\definecolor{excol}{RGB}{0,135,90}
\definecolor{attncol}{RGB}{200,40,55}
\definecolor{methcol}{RGB}{90,90,100}
\definecolor{oxcol}{RGB}{200,40,55}
\definecolor{redcol}{RGB}{0,110,180}
\definecolor{formulacol}{RGB}{198,93,9}

\newcommand{\NO}{\ensuremath{\mathrm{N.O.}}}
\newcommand{\Ered}{\ensuremath{E_{\mathrm{r}}^{\circ}}}

\tcbset{boxbase/.style={enhanced,breakable,boxrule=0.9pt,arc=2.5pt,
   left=3mm,right=3mm,top=2.3mm,bottom=2.3mm,fonttitle=\bfseries,coltitle=white,
   toptitle=0.6mm,bottomtitle=0.6mm}}
\newtcolorbox{defi}[1][]{boxbase,colback=defcol!5,colframe=defcol,title={\textsc{Définition}\ifx\relax#1\relax\relax\else\textmd{ --- \itshape #1}\fi}}
\newtcolorbox{regle}[1][]{boxbase,colback=thmcol!6,colframe=thmcol,title={\textsc{Règle}\ifx\relax#1\relax\relax\else\textmd{ --- \itshape #1}\fi}}
\newtcolorbox{exemple}[1][]{boxbase,colback=excol!5,colframe=excol,title={\textsc{Exemple}\ifx\relax#1\relax\relax\else\textmd{ : \itshape #1}\fi}}
\newtcolorbox{attention}[1][]{boxbase,colback=attncol!6,colframe=attncol,title={\textsc{Attention}\ifx\relax#1\relax\relax\else\textmd{ : \itshape #1}\fi}}

\pagestyle{fancy}\fancyhf{}
\fancyhead[L]{\small\textcolor{primarydark}{\textbf{Fiche 11} — Tuto Thème 4}}
\fancyhead[R]{\small\textcolor{primarydark}{\textsc{Potentiels standard}}}
\fancyfoot[C]{\small\thepage}
\renewcommand{\headrulewidth}{0.8pt}
\renewcommand{\headrule}{\hbox to\headwidth{\color{primary}\leaders\hrule height \headrulewidth\hfill}}
\titleformat{\section}{\large\bfseries\color{primarydark}}{\thesection}{0.6em}{}

\begin{document}
\begin{center}
{\LARGE\bfseries\color{primarydark}Thème 4 — Potentiels standard et prévision des réactions}\\[3pt]
{\color{primary}\rule{10.5cm}{1.2pt}}\\[4pt]
{\small\itshape À lire avant les exercices — comment savoir si une réaction redox aura lieu.}
\end{center}

\vspace{-2pt}
\section*{1. Le potentiel standard de réduction}
\begin{defi}[potentiel standard \Ered]
À chaque couple $\mathrm{Ox}/\mathrm{Red}$ on associe un \textbf{potentiel standard de réduction}
\Ered\ (en volts), mesuré dans les conditions standard (\SI{1}{mol\per\litre}, \SI{1}{atm},
\SI{25}{\celsius}) par rapport à l'électrode standard à hydrogène (\ce{H+}/\ce{H2}, \Ered$=0$\,V).
\end{defi}

\begin{regle}[ce que dit la valeur de \Ered]
\begin{itemize}[leftmargin=1.5em,topsep=1pt,itemsep=1.5pt]
  \item Plus \Ered\ est \textbf{élevé}, plus la forme \textbf{Ox} est un \textbf{\textcolor{oxcol}{oxydant fort}}.
  \item Plus \Ered\ est \textbf{faible (négatif)}, plus la forme \textbf{Red} est un \textbf{\textcolor{redcol}{réducteur fort}}.
\end{itemize}
\end{regle}

\section*{2. L'échelle des potentiels et la règle du « gamma »}
On classe les couples verticalement, du \Ered\ le plus élevé (en haut) au plus bas (en bas).

\begin{center}
\begin{tikzpicture}[scale=1]
  \draw[->,>=Stealth,primarydark,line width=1pt] (0,-2.6) -- (0,0.9);
  \node[primarydark,rotate=90,font=\footnotesize] at (-0.55,-1) {\Ered\ croissant};
  \foreach \y/\ox/\red/\v in {0.2/\ce{Ag+}/\ce{Ag}/+0{,}80, -1.0/\ce{Cu^{2+}}/\ce{Cu}/+0{,}34, -2.2/\ce{Zn^{2+}}/\ce{Zn}/-0{,}76}{
    \draw[formulacol,line width=1pt] (0.3,\y)--(3.7,\y);
    \node[oxcol,font=\small\bfseries] at (1.2,\y) {\ox};
    \node[redcol,font=\small\bfseries] at (3.1,\y) {\red};
    \node[primarydark,font=\scriptsize,left] at (0.2,\y) {\v};}
  \node[oxcol,font=\scriptsize\bfseries] at (1.2,0.7) {Oxydants};
  \node[redcol,font=\scriptsize\bfseries] at (3.1,0.7) {Réducteurs};
  % gamma : oxydant du haut réagit avec réducteur du bas
  \draw[->,>=Stealth,excol,line width=1.1pt] (1.2,0.0) -- (1.2,-2.0) -- (2.9,-2.1);
  \node[excol,font=\scriptsize,align=left] at (5.9,-1.0) {L'oxydant \textbf{en haut}\\ oxyde le réducteur \textbf{en bas} :\\ ici \ce{Ag+} oxyde \ce{Zn}.};
\end{tikzpicture}
\end{center}

\begin{regle}[prévoir une réaction spontanée]
Une réaction redox est \textbf{spontanée} (conditions standard) si l'\textbf{oxydant du couple de
\Ered\ le plus élevé} réagit avec le \textbf{réducteur du couple de \Ered\ le plus faible}.
De façon quantitative :
\[
\Delta E^{\circ} = \Ered(\text{oxydant}) - \Ered(\text{réducteur}) > 0 \ \text{V}.
\]
La forme dessinée par la flèche « $\gamma$ » (de l'Ox du haut vers le Red du bas) donne le sens
spontané.
\end{regle}

\section*{3. Exemples de prévision}
\begin{exemple}[\ce{Ag+} oxyde-t-il le cuivre ?]
$\Ered(\ce{Ag+}/\ce{Ag})=+0{,}80 > \Ered(\ce{Cu^{2+}}/\ce{Cu})=+0{,}34$. \ce{Ag+} (couple haut) oxyde
donc \ce{Cu} (réducteur du couple bas) : $\ce{2 Ag+ + Cu -> 2 Ag + Cu^{2+}}$.
\end{exemple}

\begin{exemple}[le fer dans l'acide chlorhydrique]
$\Ered(\ce{H+}/\ce{H2})=0{,}00 > \Ered(\ce{Fe^{2+}}/\ce{Fe})=-0{,}44$. Le proton \ce{H+} oxyde le fer :
$\ce{Fe + 2 H+ -> Fe^{2+} + H2 ^}$ (dégagement de \ce{H2}). Le fer s'oxyde en \ce{Fe^{2+}} (et non
\ce{Fe^{3+}}, car \ce{H+} ne peut pas franchir le couple \ce{Fe^{3+}}/\ce{Fe^{2+}} situé plus haut).
\end{exemple}

\begin{attention}[le sens des comparaisons]
On compare \emph{toujours} deux couples. Un oxydant ne peut oxyder un réducteur que si son couple
est \textbf{plus haut} (\Ered\ plus grand). Si $\Delta E^{\circ}<0$, la réaction proposée \textbf{n'a
pas lieu} (c'est la réaction \emph{inverse} qui est spontanée).
\end{attention}
\end{document}
